\section{平移不变体系的二次量子化哈密顿量}

Gross 第 5 章把“平面波 + 动量守恒相互作用”写成标准二次量子化形式，是 jellium、Hubbard（动量空间）与后续图技术的共同出发点。

\subsection{动能}

单粒子本征态取平面波（箱归一化 $L^d$），
\begin{equation}
  H_0
  =\sum_{\mathbf{k}\sigma}
  \varepsilon_{\mathbf{k}}
  c_{\mathbf{k}\sigma}^\dagger c_{\mathbf{k}\sigma},
  \qquad
  \varepsilon_{\mathbf{k}}=\frac{\hbar^2k^2}{2m}.
\end{equation}

\subsection{二体相互作用}

动量空间矩阵元只依赖动量转移时，
\begin{equation}
\begin{aligned}
  V
  &=\frac{1}{2L^d}
  \sum_{\mathbf{k}\mathbf{p}\mathbf{q}}
  \sum_{\sigma\sigma'}
  v_{\mathbf{q}}
  c_{\mathbf{k}+\mathbf{q},\sigma}^\dagger
  c_{\mathbf{p}-\mathbf{q},\sigma'}^\dagger
  c_{\mathbf{p}\sigma'}
  c_{\mathbf{k}\sigma}.
\end{aligned}
\label{eq:V_transl}
\end{equation}
库仑 $v_{\mathbf{q}}=4\pi e^2/q^2$（三维，$q\neq0$）；jellium 中 $q=0$ 被均匀正电背景抵消，见 \eqref{eq:jelliumH}。接触相互作用 $v_{\mathbf{q}}=U$ 即 Hubbard 的动量形式。

\subsection{为何重要}

\begin{enumerate}
  \item 总动量守恒：真空图与自能图的外线动量规则简单；
  \item $G=G(\mathbf{k},\omega)$、$\Pi=\Pi(\mathbf{q},\omega)$ 只依赖少数变量；
  \item 粒子--空穴激发、等离激元、Cooper 对都在此运动学下分类。
\end{enumerate}
后文 jellium、GMB、RPA 均默认 \eqref{eq:V_transl}。
